% Preparación separada del 30-09-2026; no modifica manuscritos publicados.
% Propietarios: VI_ES/sections/04_atlas_familias.tex:140--163;
% md/04d_ocupacion_estadistica.tex:305--319;
% md/06e_color_qutrit_gauge_finito.tex:1--194.
% Carta qutrit y forma compacta recibidas con sus hipótesis explícitas.
% Inserción: después de 00_antecedentes_geometricos_materiales y antes de 10.
\section{Résidu chromatique et représentation interne de l’action conjointe}
\label{hmtcuatro:color:seccion}

La chaîne APP--TRIT--TPK conserve l’état enrichi, ses routes, orientations, feuillets et mémoire au sein de la structure discrète conjointe du continu. Le résidu lu ci-après est une lecture finie de cet état ; il ne le remplace pas et ne réinitialise pas son histoire. Les coordonnées de phase, y compris l’unité complexe et la coordonnée régionale \(\pi\), sont des sorties HMT reçues avant cette réalisation algébrique. La représentation est construite sur le support résiduel avec la carte, l’orientation et le caractère de phase déclarés. Aucun groupe de jauge n’est inféré du seul nombre d’étiquettes.

\subsection{Support résiduel et orientation conservée}
\label{hmtcuatro:color:residuo}

Dans la copie positionnelle d’APP, soient
\[
 I_9=\{1,\ldots,9\},\qquad
 \mathcal C=I_9\times I_9,\qquad
 \mathcal C_q=\{(r,c)\in\mathcal C:r\text{ est pair}\}.
\]
L’observable résiduel et l’inversion cellulaire sont
\[
 \kappa_{\rm col}(r,c)=r-c\pmod3,\qquad
 \rho_{\rm c}(r,c)=(10-r,10-c).
\]
Le symbole \(\kappa_{\rm col}\) ne désigne pas le coefficient gravitationnel de l’action. Le relèvement de cet observable aux sections persistantes s’effectue par la projection sur l’atlas et conserve les autres registres ; il ne sélectionne pas une cellule par sa masse évaluée.

\begin{proposition}[Fibres ternaires et bilan chromatique]
\label{hmtcuatro:color:balance} La fonctionnelle \(\kappa_3:\mathbb F_3^2\to\mathbb F_3\), \(\kappa_3(r,c)=r-c\), possède trois fibres affines de cardinal trois. Sur \(\mathcal C_q\), les trois résidus ont pour cardinaux \(12,12,12\). L’inversion cellulaire change le signe du résidu.
\end{proposition}
\begin{proof}
Pour chaque \(j\in\mathbb F_3\), la fibre est \(\{(c+j,c):c\in\mathbb F_3\}\), un translaté du noyau \(\{(c,c)\}\). Dans \(\mathcal C_q\), il existe quatre valeurs paires de \(r\) ; pour chacune, \(c=1,\ldots,9\) parcourt trois fois chaque classe modulo trois. Chaque résidu possède \(4\cdot3=12\) représentants. Enfin,
\[
 \kappa_{\rm col}(\rho_{\rm c}(r,c))
 =(10-r)-(10-c)=-(r-c)\pmod3.
\]
La classe zéro reste fixe et les classes un et deux sont échangées.
\end{proof}

Ce bilan n’identifie pas saveur, profondeur et couleur. Il n’affirme pas non plus que les trois droites qui seront construites sont trois états complets : les routes et la mémoire qui partagent une lecture résiduelle demeurent dans la fibre enrichie.

\subsection{Carte qutrit, paire de Weyl et algèbre de couleur}
\label{hmtcuatro:color:weyl}

On fixe la carte
\[
 V_{\rm col}=\mathbb C[\mathbb F_3]\simeq\mathbb C^3,
 \qquad (e_0,e_1,e_2),\qquad
 \omega_3=e^{2\pi i/3},
\]
avec une base delta orthonormée, l’orientation \(j\mapsto j+1\) et le caractère \(j\mapsto\omega_3^j\). Les opérateurs sont
\[
 Xe_j=e_{j+1\bmod3},\qquad Ze_j=\omega_3^j e_j,
 \qquad X^3=Z^3=I,\quad ZX=\omega_3XZ.
\]
La dernière égalité se vérifie en appliquant les deux membres à chaque \(e_j\). L’application de l’atlas vers cette carte est
\[
 (r,c)\longmapsto\mathbb C e_{\kappa_{\rm col}(r,c)}.
\]
Ainsi, \(X\) permute les droites et \(Z\) enregistre leurs phases. Changer la carte affine et transporter conjointement base, orientation et caractère conjugue les opérateurs par la matrice unitaire de changement de base correspondante. Cela ne revient pas à changer une seule étiquette en laissant les autres données fixes.

\begin{theorem}[Algèbre matricielle et forme réelle compacte]
\label{hmtcuatro:color:algebra} Les neuf opérateurs \(W_{ab}=X^aZ^b\), \((a,b)\in\mathbb F_3^2\), forment une base orthogonale de \(\operatorname{End}_{\mathbb C}(V_{\rm col})\). Les huit indices non nuls engendrent linéairement \(\mathfrak{sl}_3(\mathbb C)\), dont la forme réelle de matrices antihermitiennes est \(\mathfrak{su}(3)\).
\end{theorem}
\begin{proof}
Si \(a\ne0\), la matrice \(X^aZ^b\) n’a pas d’entrées diagonales. Si \(a=0\), sa trace est \(1+\omega_3^b+\omega_3^{2b}\), qui vaut zéro pour \(b\ne0\) et trois pour \(b=0\). La relation de Weyl donne alors
\[
 \operatorname{tr}(W_{ab}^{\dagger}W_{cd})
 =3\delta_{ac}\delta_{bd}.
\]
L’indépendance et la dimension neuf prouvent la première assertion ; les huit autres forment une base du sous-espace de trace nulle. De plus,
\[
 W_{ab}^{\dagger}=\omega_3^{ab}W_{-a,-b}.
\]
Les huit indices non nuls forment quatre paires adjointes. En prenant \(\mathcal R=\{(1,0),(0,1),(1,1),(1,2)\}\), les matrices
\[
 A_u=W_u-W_u^{\dagger},\qquad
 B_u=i(W_u+W_u^{\dagger}),\qquad u\in\mathcal R,
\]
sont antihermitiennes et de trace nulle. Leur indépendance réelle découle de l’indépendance complexe des huit \(W_u,W_{-u}\) : dans chaque paire, les coefficients d’une combinaison réelle nulle imposent l’annulation des deux coefficients réels. Elles forment une base réelle de dimension huit. Par conséquent,
\[
 \{A\in\mathfrak{sl}_3(\mathbb C):-A^{\dagger}=A\}
 =\mathfrak{su}(3).
\]
\end{proof}

La forme hermitienne et l’orientation complexe fixées déterminent le groupe de réalisation
\[
 SU(V_{\rm col})=\{U:U^{\dagger}U=I,\ \det U=1\}.
\]
Son algèbre tangente est exactement la précédente : différencier les deux équations à l’identité donne \(A^{\dagger}=-A\) et \(\operatorname{tr}A=0\) ; réciproquement, \(e^{tA}\) satisfait les deux. Tout élément du groupe s’écrit comme l’exponentielle d’une matrice de cette algèbre. En effet, diagonalisons unitairement \(U\) avec les valeurs propres \(e^{i\theta_j}\) ; comme le déterminant vaut un, \(\sum_j\theta_j=2\pi n\). Soustraire \(2\pi n\) à l’un des angles ne modifie pas \(U\) et produit un logarithme antihermitien de trace nulle. Cette réalisation continue ne s’identifie pas au groupe fini engendré uniquement par \(X,Z\). Elle conserve la forme compacte et n’introduit aucune valeur pour le couplage fort.

\subsection{Transport chromatique et facteurs internes indépendants}
\label{hmtcuatro:color:transporte}

Dans une cellularisation finie orientée compatible avec les routes, on attribue \(U_e\in SU(V_{\rm col})\) et \(U_{\bar e}=U_e^{-1}\). Pour \(g_v\in SU(V_{\rm col})\),
\[
 U_e^g=g_{t(e)}U_eg_{s(e)}^{-1},\qquad
 \psi_v^g=g_v\psi_v.
\]
Par substitution,
\[
 U_e^g\psi_{s(e)}^g-\psi_{t(e)}^g
 =g_{t(e)}(U_e\psi_{s(e)}-\psi_{t(e)}).
\]
Le signe opposé, utilisé dans \(D_T\), possède la même covariance. La transformation d’un produit orienté autour d’une face est une conjugaison par la jauge du sommet de base : les facteurs des sommets intermédiaires s’annulent successivement. En particulier, pour \(\beta\ge0\),
\[
 S_{\rm col}^{(\mathcal K,\beta)}[U]
 =\frac\beta3\sum_f(3-\operatorname{Re}\operatorname{tr}U_f)
\]
est invariante et positive ou nulle. L’invariance utilise la trace cyclique ; la positivité résulte de \(\operatorname{Re}\operatorname{tr}U_f\le3\). Le paramètre et la cellularisation conservent leur statut de données de cette réalisation. Cette preuve ne les détermine pas par comparaison avec des valeurs chromodynamiques.

La matière de l’action conjointe utilise ce facteur, non une autre partition de l’atlas. Pour les quarks gauches, sa fibre interne possède les facteurs \(V_{\rm col}\otimes\mathbb C^2\otimes F\), où \(F\) conserve les familles ; pour chaque type droit, on utilise \(V_{\rm col}\otimes F\). Le facteur de couleur est trivial pour les leptons. Le facteur spinoriel est conservé séparément. Pour des matrices de couleur \(A\) et des matrices faibles \(B\), on a
\[
 [A\otimes I,I\otimes B]=0,
\]
car les deux produits valent \(A\otimes B\). Les caractères d’hypercharge agissent comme des scalaires sur chaque multiplet et commutent également. Cette identité permet de composer les connexions internes et celle de spin dans la même dérivée sans confondre couleur, chiralité et saveur.

Les applications de Yukawa de l’action ultérieure sont l’identité sur le facteur chromatique et des applications sur les facteurs Higgs--famille. Par conséquent, \(Y_q(H)(g_{\rm col}\otimes I)
=(g_{\rm col}\otimes I)Y_q(H)\). La compatibilité faible et d’hypercharge y est démontrée avec leurs colonnes et leurs poids explicites ; elle ne résulte pas du recensement \(12+12+12\). On conserve ainsi la chaîne résidu--qutrit--algèbre--transport requise par l’action commune, sans affirmer dans ces fondements une limite quantique de Yang--Mills ni remplacer ses documents sources.
